\subsection{The transported shock budget gives a graph-uniform terminal horizon}
\label{sec:complete-gate-horizon}

The quartic score reserve is not itself a convergence potential.  For the
same transported-center recurrence, however, the total energy before
subtracting the consumed reserve has a graph-uniform cap.  This gives a
direct terminal horizon for the complete gate.

\begin{theorem}[Graph-uniform soft terminal horizon for the complete gate]
\label{thm:complete-gate-soft-horizon}
Retain the finite-graph, exact-real, point-seed recurrence, frozen
zero-padding, exact transported centers, original safe envelope, and named
complete gate of \cref{thm:original-score-projected-quartic}.  Thus
\begin{equation*}
 0<q<1,
 \qquad \alpha=q^2,
 \qquad \rho=\tau=\frac q5,
\end{equation*}
the initial face is $U_0=\{s\}$, its restricted optimum is interior, every
fixed-face step is followed immediately by the complete gate, and admissions
contain exactly all genuinely violating boundary rows.  Define
\begin{align}
 E_{\rm gate}(q)
 &:=\frac{q^{12}}{50(1+q^2)^2},
 \label{eq:complete-gate-energy-threshold}\\
 K_q
 &:=\left\lceil
 \frac{\log\!\left(50(1+q^2)^2q^{-10}\right)}
      {-\log(1-q)}
 \right\rceil.
 \label{eq:complete-gate-hold-bound}
\end{align}
Then every visited face undergoes at most $K_q$ consecutive fixed-face steps
before the gate either admits a nonempty batch or returns the global
one-sided certificate.  In particular, if $N_{\rm exp}$ is the number of
nonempty admission events, $T$ the number of fixed-face steps, and
\begin{equation*}
 \mathfrak V_T:=\sum_{t=0}^{T-1}\vol(U_t),
\end{equation*}
then
\begin{align}
 N_{\rm exp}
 &\leq \abs{S_\rho\setminus U_0}\leq\frac5q,
 \label{eq:complete-gate-admission-count}\\
 T
 &\leq(N_{\rm exp}+1)K_q
 \leq\left(\frac5q+1\right)K_q
 =\bigO\!\left(q^{-2}\log\frac1q\right),
 \label{eq:complete-gate-total-steps}\\
 \mathfrak V_T
 &\leq\frac5qT
 =\bigO\!\left(q^{-3}\log\frac1q\right)
 \qquad(q\downarrow0).
 \label{eq:complete-gate-swept-volume}
\end{align}
The execution therefore terminates after finitely many steps with a safe
envelope $\ell_T$ satisfying the note-scoped PPR guarantee
\begin{equation}
 \norm{D^{-1/2}(x^0-\ell_T)}_\infty\leq\rho+\tau=\frac{2q}{5}.
 \label{eq:complete-gate-terminal-error}
\end{equation}
These are graph-uniform convergence, fixed-face chronology, and swept-volume
bounds for this exact recurrence.  They are not the desired product-scale
work theorem.
\end{theorem}

\begin{proof}
Write the recurrence in normalized coordinates, as in
\cref{prop:consumed-energy-structural-solvency}.  Thus
$B_U=D_UH_U$ is symmetric positive definite,
\begin{equation*}
 \Phi_U(z)=\frac12z^\top B_Uz-h_U^\top D_Uz,
 \qquad
 h_i=q^2\left(\frac{\mathbf1_{\{i=s\}}}{d_i}-\frac q5\right),
\end{equation*}
and $E_U(z,v)$ is given by \cref{eq:consumed-reserve-energy}.  Let $D_k$ be
the sum of all Schur drops through the current face.  Exact transport and the
restricted-optimum telescope give the identity
\begin{align}
 E_0+D_k
 &= -\Phi_{U_k}(z_{U_k}^\star)
    +\frac{q^2}{2}\norm{z_{U_0}^\star}_{D_{U_0}}^2.
 \label{eq:complete-gate-total-energy-identity}
\end{align}
Indeed, the objective-gap part of the actual zero-start energy cancels the
first restricted optimum against the unweighted Schur telescope.

The declared interiority propagates from the initial singleton.  Write
$\bar\ell:=D^{-1/2}\ell$ for the normalized safe envelope.  If $B$ is one
admitted batch, envelope safety and the nonpositive cut block give
\begin{equation*}
 r_B(z_U^\star)\leq r_B(\bar\ell)<-q^2\tau\one.
\end{equation*}
The Schur complement on $B$ is a positive definite $M$-matrix.  Hence its
new-coordinate correction is $-S_B^{-1}r_B(z_U^\star)>0$, and the induced
old-coordinate correction is nonnegative.  Thus the enlarged restricted
optimum is again interior, including for a simultaneous complete batch.

This total energy is graph-uniformly bounded.  In the $D$ inner product,
$B\succeq q^2D$.  Since the point-seed load has at most one positive entry,
\begin{equation*}
 \norm{h_+}_D^2
 =d_sh_s^2
 \leq\frac{q^4}{d_s}\leq q^4.
\end{equation*}
Completing the square, and then comparing a restricted optimum with the full
obstacle optimum, yields
\begin{equation*}
 -\Phi_{U_k}(z_{U_k}^\star)
 \leq-\Phi(z_\rho^\star)
 \leq\frac{\norm{h_+}_D^2}{2q^2}
 \leq\frac{q^2}{2}.
\end{equation*}
If $a=(1+q^2)/2$, then the initial singleton optimum is
$z_{U_0,s}^\star=h_s/a$, so
\begin{equation*}
 \frac{q^2}{2}\norm{z_{U_0}^\star}_{D_{U_0}}^2
 \leq\frac{q^6}{2a^2}
 =\frac{2q^6}{(1+q^2)^2}
 \leq\frac{q^2}{2}.
\end{equation*}
The last inequality is equivalent to
$(3q^2+1)(q^2-1)\leq0$.  Therefore
\begin{equation}
 E_k\leq E_0+D_k\leq q^2
 \label{eq:complete-gate-graph-uniform-energy-cap}
\end{equation}
at every checkpoint.  The first inequality is exactly nonnegativity of the
consumed reserve; it retains rather than discards every admission shock.

We next turn small energy into a gate decision.  On a current interior face
put
\begin{equation*}
 e:=\norm{z-z_U^\star}_\infty.
\end{equation*}
The objective-gap part of $E$ and $B_U\succeq q^2D_U$ imply
\begin{equation}
 e\leq\frac{\sqrt{2E}}q.
 \label{eq:complete-gate-energy-to-coordinate}
\end{equation}
Every row of $H_U$ has diagonal magnitude $(1+q^2)/2$ and off-diagonal
absolute row sum at most $(1-q^2)/2$, hence
$\norm{H_U}_{\infty\to\infty}\leq1$.  Interiority gives
$r_U(z)=H_U(z-z_U^\star)$, so the original correction and envelope obey
\begin{align}
 \delta(z)&\leq\frac e{q^2},
&\norm{\bar\ell-z_U^\star}_\infty
 &\leq(1+q^{-2})e.
 \label{eq:complete-gate-envelope-distance}
\end{align}
Here the second inequality also uses nonexpansiveness of the coordinatewise
positive part and $z_U^\star\geq0$; it does not assume inactive projection.
If $E\leq E_{\rm gate}(q)$, then
\cref{eq:complete-gate-energy-to-coordinate,eq:complete-gate-envelope-distance}
give
\begin{equation*}
 \norm{\bar\ell-z_U^\star}_\infty
 \leq\frac{q^3}{5}=q^2\tau.
\end{equation*}
Applying the same row-sum bound to $H_U(\bar\ell-z_U^\star)$ shows that every
active row is at least $-q^2\tau$.  The complete gate explicitly checks every
boundary row.  Every nonseed vertex outside $U\cup N(U)$ has normalized KKT
value $q^2\rho>0$, while the seed belongs to $U_0$.  Thus the gate either
finds and admits a nonempty genuinely violating boundary batch or all rows
satisfy the global certificate.

At entry to any face, \cref{eq:complete-gate-graph-uniform-energy-cap} gives
$E\leq q^2$.  After $k$ fixed-face steps,
\cref{eq:fixed-face-center-contraction} gives
$E\leq(1-q)^kq^2$.  The definition of $K_q$ makes this no larger than
$E_{\rm gate}(q)$, proving the hold bound.

Finally, safe admission keeps every active face inside $S_\rho$, whose degree
volume is at most $1/\rho=5/q$.  Every nonempty admission adds at least one
new support vertex.  More precisely, the chronology counts the initial
singleton face once and creates one subsequent face only after each nonempty
complete-gate batch.  The admitted batches are pairwise disjoint subsets of
$S_\rho\setminus U_0$; empty decisions and forbidden closure rows create no
faces.  This proves \cref{eq:complete-gate-admission-count} and gives exactly
$N_{\rm exp}+1$ visited faces.  Summing the per-face hold and volume bounds proves
\cref{eq:complete-gate-total-steps,eq:complete-gate-swept-volume}.  Since
$-\log(1-q)\geq q$, $K_q=O(q^{-1}\log(1/q))$ as $q\downarrow0$.  Finite
termination follows.  At the terminal certificate, the maximum-principle
bound and the RPPR regularization bias give the explicit semantic conversion
\begin{align*}
 \norm{D^{-1/2}(x^0-\ell_T)}_\infty
 &\leq
 \norm{D^{-1/2}(x^0-x_\rho^\star)}_\infty
 +\norm{D^{-1/2}(x_\rho^\star-\ell_T)}_\infty\\
 &\leq\rho+\tau=\frac{2q}{5},
\end{align*}
by \cref{thm:rppr-support,lem:one-sided-certificate}.  This proves
\cref{eq:complete-gate-terminal-error}.
\end{proof}

\begin{corollary}[A charged exact-real fallback, not product-scale work]
\label{cor:complete-gate-dense-fallback}
Under the hypotheses of \cref{thm:complete-gate-soft-horizon}, suppose the
implementation caches every admitted adjacency row, evaluates each fixed-face
step and complete boundary gate by scanning the current active incidences,
and maintains exact restricted optima and transported displacements with a
dense incremental Cholesky factor.  In the unit-cost exact-real arithmetic
model, active recurrence and gate work is
\begin{equation*}
 O(\mathfrak V_T)=O(q^{-3}\log(1/q)),
\end{equation*}
cumulative factor updates and triangular solves cost $O(n_{\rm fin}^3)$,
and persistent factor storage costs $O(n_{\rm fin}^2)$.  Since
$n_{\rm fin}\leq\vol(U_T)\leq5/q$, the total charged arithmetic is
\begin{equation}
 O(q^{-3}\log(1/q)),
 \label{eq:complete-gate-dense-work}
\end{equation}
with $O(q^{-2})$ persistent cells, $O(q^{-1})$ one-time admitted-row
incidence exposure, and $O(q^{-1})$ terminal output cells.

This fallback charges every optimum/transport query but assumes exact-real
unit-cost arithmetic and a constant-cost degree query for boundary rows.  It
gives no finite-precision, bit-complexity, numerical-stability, sparse-fill,
or intermediate-emission guarantee.  In particular,
\cref{eq:complete-gate-dense-work} is a factor
$\widetilde O(q^{-1})$ above the desired constant-ratio product target
$\widetilde O(1/(\rho q))=\widetilde O(q^{-2})$.
\end{corollary}

\begin{proof}
A fixed-face recurrence and projection scan the cached active incidences and
write the new primal and center vectors.  The active residual maximum and
complete boundary residual accumulation use another constant number of such
scans; boundary degree queries are charged explicitly.  Thus all repeated
old-face adjacency reads, arithmetic, envelope materialization, and state
writes cost $O(\vol(U_t))$ per fixed-face step and
$O(\mathfrak V_T)$ in total.  Reading every newly admitted adjacency row once
costs its ambient degree, so safe support gives total one-time exposure at
most $5/q$.  For a final active dimension
$n_{\rm fin}$, block Cholesky updates over any partition into admission
batches cost $O(n_{\rm fin}^3)$ in total.  One pair of triangular solves per
visited face also sums to $O(n_{\rm fin}^3)$ because there are at most
$n_{\rm fin}$ nonempty admissions and every solve costs
$O(n_{\rm fin}^2)$.  Materializing and writing each transported displacement
costs $O(n_{\rm fin})$ per face and is dominated by those solves.  The factor
has $O(n_{\rm fin}^2)$ cells.  One final cached-row certificate scan and the
$O(n_{\rm fin})$-cell output are charged separately.  Substitute
\cref{eq:complete-gate-swept-volume} and
$n_{\rm fin}\leq5/q$.
\end{proof}
